% BÀI TOÁN HÌNH HỌC - BAI_HINH_55 (TRÍCH ĐỀ 55)
\begin{ex}
\begin{ex}
Cho đường tròn $(O; R)$ và điểm $A$ nằm ngoài đường tròn sao cho $OA > 2R$. Từ $A$ kẻ hai tiếp tuyến $AB, AC$ với đường tròn $(O)$ ($B, C$ là các tiếp điểm). Gọi $H$ là giao điểm của $OA$ và $BC$. Kẻ cát tuyến $ADE$ với đường tròn $(O)$ ($D$ nằm giữa $A$ và $E$, tia $AE$ nằm giữa hai tia $AO$ và $AC$).
\begin{enumerate}[label=\textbf{\alph*)}]
    \item Chứng minh tứ giác $ABOC$ nội tiếp và $OA \perp BC$ tại $H$.
    \item Chứng minh $AH \cdot AO = AD \cdot AE = AB^2$.
    \item Đường thẳng qua $D$ song song với $BE$ cắt $AB, BC$ lần lượt tại $I$ và $K$. Chứng minh $D$ là trung điểm của $IK$.
\end{enumerate}
\loigiai{
{
    \begin{itemize}
        \item[a)] Tứ giác $ABOC$ có $\widehat{ABO} = \widehat{ACO} = 90^\circ \Rightarrow \widehat{ABO} + \widehat{ACO} = 180^\circ \Rightarrow ABOC$ nội tiếp đường tròn đường kính $AO$. Vì $AB = AC$ và $OB = OC = R$ nên $OA$ là trung trực của $BC \Rightarrow OA \perp BC$ tại $H$.
        \item[b)] Xét $\triangle ABD$ và $\triangle AEB$ có $\widehat{BAE}$ chung, $\widehat{ABD} = \widehat{AEB}$ (góc tạo bởi tia tiếp tuyến và dây cung và góc nội tiếp cùng chắn $\stackrel{\frown}{BD}$). Suy ra $\triangle ABD \sim \triangle AEB$ (g-g) $\Rightarrow \dfrac{AB}{AE} = \dfrac{AD}{AB} \Rightarrow AB^2 = AD \cdot AE$. Áp dụng hệ thức lượng vào $\triangle ABO$ vuông tại $B$, đường cao $BH$: $AB^2 = AH \cdot AO$. Suy ra $AH \cdot AO = AD \cdot AE = AB^2$.
        \item[c)] Sử dụng định lý Thales cho các đường song song và tính chất chùm điều hòa / tứ giác nội tiếp suy ra $D$ là trung điểm của $IK$.
    \end{itemize}
}
\end{ex
}
\end{ex}
